\documentclass[11pt,a4paper]{article}
\usepackage[T1]{fontenc}
\usepackage[utf8]{inputenc}
\usepackage[spanish,es-nodecimaldot,es-noquoting,es-lcroman]{babel}
\spanishplainpercent
\usepackage{lmodern}
\usepackage[margin=2.3cm,headheight=15pt]{geometry}
\usepackage{amsmath,amssymb,bm}
\usepackage{booktabs,array,longtable}
\usepackage{tikz}
\usetikzlibrary{arrows.meta,positioning,calc,matrix,decorations.pathreplacing}
\usepackage[most]{tcolorbox}
\usepackage{listings}
\usepackage{microtype}
\usepackage{fancyhdr}
\usepackage{hyperref}
\usepackage{bookmark}
\hypersetup{colorlinks=true,linkcolor=blue!45!black,urlcolor=blue!50!black,
 pdftitle={BEMCilindro explicado linea por linea},
 pdfauthor={Documento de estudio del notebook bigardi.ipynb}}
\setlength{\parindent}{0pt}
\setlength{\parskip}{0.45em}
\setcounter{tocdepth}{2}
\numberwithin{equation}{section}
\pagestyle{fancy}
\fancyhf{}
\fancyhead[L]{\small \texttt{BEMCilindro}: explicaci\'on paso a paso}
\fancyhead[R]{\small Documento independiente}
\fancyfoot[C]{\thepage}
\newcommand{\vect}[1]{\bm{#1}}
\newcommand{\dd}{\,\mathrm{d}}
\newcommand{\ii}{\mathrm{i}}
\newcommand{\norm}[1]{\left\lVert#1\right\rVert}
\DeclareRobustCommand{\code}[1]{\texttt{\detokenize{#1}}}
\newtcolorbox{idea}[1][Idea clave]{colback=blue!3,colframe=blue!45!black,
 title={#1},fonttitle=\bfseries,breakable}
\newtcolorbox{cuidado}[1][No confundir]{colback=orange!4,
 colframe=orange!60!black,title={#1},fonttitle=\bfseries,breakable}
\lstset{language=Python,basicstyle=\ttfamily\small,
 keywordstyle=\color{blue!60!black},commentstyle=\color{green!35!black},
 stringstyle=\color{orange!55!black},frame=single,
 rulecolor=\color{black!20},breaklines=true,columns=fullflexible,
 keepspaces=true,showstringspaces=false,upquote=true,
 escapeinside={(*@}{@*)},numberstyle=\tiny\color{black!55},
 numbersep=8pt,xleftmargin=1em,
 postbreak=\mbox{\textcolor{gray}{$\hookrightarrow$}\space}}

\begin{document}
\begin{titlepage}
\centering
\vspace*{1.4cm}
{\Large Lectura guiada de c\'odigo cient\'ifico\par}
\vspace{1.1cm}
{\Huge\bfseries \texttt{BEMCilindro}\\[0.35cm]
explicado l\'inea por l\'inea\par}
\vspace{0.8cm}
{\Large Qu\'e hace cada instrucci\'on, qu\'e guarda cada arreglo\\
y c\'omo se construye y resuelve el sistema BEM\par}
\vspace{1.5cm}
\begin{tikzpicture}[>=Stealth,node distance=0.65cm,
 box/.style={draw,rounded corners,align=center,text width=10cm,
 minimum height=1.1cm,fill=blue!4}]
 \node[box] (a){\code{__init__}\\Preparar la geometr\'ia, la cuadratura y $H^s$};
 \node[box,below=of a] (b){\code{matrices(w)}\\Construir $H(\omega)$ y $G(\omega)$};
 \node[box,below=of b] (c){\code{resolver(w, p)}\\Aplicar la presi\'on, resolver y proyectar radialmente};
 \draw[->,thick] (a)--(b);
 \draw[->,thick] (b)--(c);
\end{tikzpicture}
\vfill
{\large Documento complementario a la gu\'ia general del BEM\par}
{\large Ejemplo comprobado: $R=6$ m, 8 elementos, $\omega=40$ rad/s\par}
\vspace{0.7cm}
{\large Noviembre de 2026\par}
\end{titlepage}

\tableofcontents
\clearpage

\section{Objetivo y manera de leer este documento}
Este documento explica exclusivamente la clase \code{BEMCilindro}
compartida por el usuario. No reemplaza la gu\'ia general ni modifica el
notebook. El c\'odigo compartido se compar\'o con la celda de la clase
en \code{estudio_bignardi/bigardi.ipynb}: sus instrucciones coinciden.

La explicaci\'on combina tres lecturas de cada operaci\'on:
\begin{enumerate}
 \item \textbf{Python:} qu\'e hace la instrucci\'on.
 \item \textbf{Arreglos:} cu\'ales son las dimensiones y los \'indices.
 \item \textbf{BEM:} qu\'e cantidad f\'isica o integral se est\'a calculando.
\end{enumerate}
Las l\'ineas citadas corresponden al listado completo de la
secci\'on~\ref{sec:listado}. Una l\'inea larga puede continuar visualmente
en otro rengl\'on, sin cambiar su n\'umero.

\begin{idea}[Lo que resuelve la clase]
Con geometr\'ia y material conocidos, se prescribe la presi\'on en todo
el borde. El BEM produce
\begin{equation}
 \boxed{H(\omega)\vect u=G(\omega)\vect t,\qquad
 \vect t=-p\vect n.}
 \label{eq:sistema}
\end{equation}
$\vect t$ es conocido y $\vect u$ es desconocido. La clase resuelve ese
sistema y devuelve el desplazamiento radial complejo en cada nodo.
\end{idea}

\subsection{Qu\'e no se volver\'a a deducir}
Se da por conocida la funci\'on de cada n\'ucleo:
$U$ relaciona fuerzas auxiliares con desplazamientos y $T$ con tracciones.
La deducci\'on completa de esos n\'ucleos pertenece a la gu\'ia general.
Aqu\'i se explica \emph{c\'omo se utilizan e integran}, con sus dimensiones,
signos e \'indices.

\subsection{Caso concreto que acompa\~na la lectura}
Se utilizar\'a
\begin{equation}
 R=6\ {\rm m},\quad N_e=8,\quad N=16,\quad
 p=100\ {\rm Pa},\quad\omega=40\ {\rm rad/s}.
\end{equation}
El material es el del notebook:
$\rho=100$ kg/m$^3$, $\mu_0=10^6$ Pa, $\nu=0.25$ y $\beta=0.05$.
Los n\'umeros de ejemplo se obtuvieron ejecutando las funciones originales,
no estim\'andolos a partir de una figura.

\clearpage
\section{El listado completo y su numeraci\'on}
\label{sec:listado}
El listado conserva las 49 l\'ineas del fragmento, incluidos los comentarios
y los separadores en blanco. Los comentarios acentuados se representan
tipogr\'aficamente en \LaTeX{}; las instrucciones Python son las mismas.

\begin{lstlisting}[numbers=left,firstnumber=1,label=lst:clase,
 basicstyle=\ttfamily\footnotesize,breakautoindent=false,breakindent=1em]
class BEMCilindro:
    def __init__(self, R, ne, m):
        self.R, self.ne, self.m = R, ne, m
        self.nodes, self.conn = malla_circulo(R, ne)
        self.N = len(self.nodes)
        # Prec(*@\'a@*)lculo de cuadratura para cada par (nodo de colocaci(*@\'o@*)n, elemento)
        self.quad = []
        Hs = np.zeros((2*self.N, 2*self.N))
        for ic, xi in enumerate(self.nodes):
            fila = []
            for cn in self.conn:
                loc = np.where(cn == ic)[0]
                eta_s = [-1.0, 0.0, 1.0][loc[0]] if len(loc) else None
                eta, wq = puntos_cuadratura(eta_s)
                Pn, dP = shape(eta), dshape(eta)
                X, dX = Pn @ self.nodes[cn], dP @ self.nodes[cn]
                J = np.linalg.norm(dX, axis=-1)
                nrm = np.stack([dX[:, 1], -dX[:, 0]], -1)/J[:, None]
                W = Pn*(J*wq)[:, None]                          # P_n(eta) J w_q
                Ts = kernel_Ts(X, xi, nrm, m["nu"])
                for a, g in enumerate(cn):
                    if g != ic:                                 # bloque diagonal: s(*@\'o@*)lido r(*@\'i@*)gido
                        Hs[2*ic:2*ic+2, 2*g:2*g+2] += np.einsum("q,qij->ij", W[:, a], Ts)
                fila.append((cn, X, nrm, W, Ts))
            self.quad.append(fila)
        for ic in range(self.N):
            Hs[2*ic:2*ic+2, 2*ic:2*ic+2] = -sum(Hs[2*ic:2*ic+2, 2*g:2*g+2] for g in range(self.N) if g != ic)
        self.Hs = Hs

    def matrices(self, w):
        N = self.N
        G = np.zeros((2*N, 2*N), complex)
        H = self.Hs.astype(complex)
        for ic, xi in enumerate(self.nodes):
            for cn, X, nrm, W, Ts in self.quad[ic]:
                U, T = kernels(X, xi, nrm, w, self.m)
                Gb = np.einsum("qa,qij->aij", W, U)
                Hb = np.einsum("qa,qij->aij", W, T - Ts)       # parte din(*@\'a@*)mica regular
                for a, g in enumerate(cn):
                    G[2*ic:2*ic+2, 2*g:2*g+2] += Gb[a]
                    H[2*ic:2*ic+2, 2*g:2*g+2] += Hb[a]
        return H, G

    def resolver(self, w, p=P):
        H, G = self.matrices(w)
        nrm = self.nodes/self.R
        t = (-p*nrm).ravel()                                   # t = -p n en todos los nodos
        u = np.linalg.solve(H, G @ t).reshape(-1, 2)
        return (u*nrm).sum(-1)                                 # u_r en cada nodo
\end{lstlisting}

\clearpage
\section{Panorama: tres m\'etodos y tres momentos distintos}
\subsection{Crear un objeto no es resolver una frecuencia}
La instrucci\'on
\begin{lstlisting}
bem = BEMCilindro(R, 8, mat)
\end{lstlisting}
crea una instancia y ejecuta autom\'aticamente \code{__init__}.
Prepara la malla, los datos de integraci\'on y $H^s$.
No se ha elegido una frecuencia de resoluci\'on en ese constructor.

Despu\'es:
\begin{lstlisting}
H, G = bem.matrices(40.0)
ur = bem.resolver(40.0)
\end{lstlisting}
La primera llamada entrega matrices, sin resolver el desplazamiento.
La segunda \emph{vuelve a construir las matrices} internamente y resuelve.
No aprovecha autom\'aticamente las variables locales \code{H} y \code{G}
de la llamada anterior.

\begin{cuidado}
Para usar normalmente la clase basta crear el objeto y llamar a
\code{resolver}. Llamar antes a \code{matrices} solo es necesario si
se quieren inspeccionar las matrices; de otro modo se duplica
innecesariamente ese ensamblaje.
\end{cuidado}

\subsection{Qu\'e permanece y qu\'e se recalcula}
\begin{center}
\begin{tabular}{p{0.29\textwidth}p{0.61\textwidth}}
\toprule
Momento & Datos principales\\
\midrule
Una vez, en \code{__init__}
& Nodos, conectividad, puntos f\'isicos de cuadratura, normales,
pesos combinados, $T^s$ y $H^s$.\\
Cada \code{matrices(w)}
& $U(\omega)$, $T(\omega)$, $G(\omega)$ y $H(\omega)$.\\
Cada \code{resolver(w,p)}
& Lo anterior, la carga, la soluci\'on cartesiana y su proyecci\'on radial.\\
\bottomrule
\end{tabular}
\end{center}

\subsection{Qu\'e significa \texttt{self}}
\code{self} es la referencia al objeto que recibe la llamada.
En \code{bem.resolver(40.0)}, Python suministra \code{bem} como
primer argumento del m\'etodo. El usuario no debe escribirlo de nuevo.

\begin{lstlisting}
self.R = R
\end{lstlisting}
guarda el radio en el objeto. En cambio:
\begin{lstlisting}
N = self.N
\end{lstlisting}
crea un nombre local dentro del m\'etodo. El atributo
\code{self.N} sigue existiendo cuando el m\'etodo termina;
el nombre local \code{N} no se convierte en un nuevo atributo.

\section{Diccionario de nombres, \'indices y dimensiones}
\label{sec:dimensiones}
\subsection{No confundir nombres parecidos}
\begin{longtable}{p{0.20\textwidth}p{0.70\textwidth}}
\toprule
Nombre & Significado\\
\midrule
\endhead
\code{ne} & N\'umero de elementos.\\
\code{self.N}, \code{N} & N\'umero de nodos, no de elementos.\\
\code{w} & Frecuencia angular $\omega$, escalar, en rad/s.\\
\code{wq} & Pesos de cuadratura; no son frecuencias.\\
\code{W} & Funciones de forma por jacobiano por peso de cuadratura.\\
\code{P} & Presi\'on global que fija el valor predeterminado del m\'etodo.\\
\code{p} & Presi\'on suministrada a \code{resolver}.\\
\code{Pn} & Valores de funciones de forma, no presiones.\\
\code{xi} & Coordenadas del nodo fuente/de colocaci\'on.\\
\code{X} & Coordenadas de todos los puntos de cuadratura de un elemento.\\
\code{Ts} & Valores del n\'ucleo est\'atico en puntos de cuadratura.\\
\code{Hs} & Matriz est\'atica global ensamblada.\\
\code{Hb}, \code{Gb} & Tres bloques locales, a\'un no matrices globales.\\
\bottomrule
\end{longtable}

\subsection{Los cuatro tipos de \'indice}
Utilizaremos:
\begin{equation}
 c=\text{nodo de colocaci\'on},\quad
 e=\text{elemento},\quad
 a=\text{nodo local},\quad g=\text{nodo global}.
\end{equation}
En el c\'odigo, $c$ se llama \code{ic}.
El bucle de elementos no guarda expl\'icitamente $e$:
usa directamente su conectividad \code{cn}.
Adem\'as:
\begin{equation}
 q=\text{punto de cuadratura},\qquad i,j=0,1
   =\text{componentes espaciales}.
\end{equation}
Los \'indices $i,j$ de los n\'ucleos no son n\'umeros de nodo.
En la convenci\'on del notebook, el primero identifica la direcci\'on
de la fuerza auxiliar y la ecuaci\'on; el segundo, la componente
de campo que se integra.

\subsection{Tabla de formas}
Si $n_q$ es el n\'umero de puntos usados en el par actual:
\begin{longtable}{lll}
\toprule
Arreglo & Forma & En el ejemplo habitual\\
\midrule
\endhead
\code{nodes} & $(N,2)$ & $(16,2)$\\
\code{conn} & $(N_e,3)$ & $(8,3)$\\
\code{cn} & $(3,)$ & Tres \'indices enteros\\
\code{eta}, \code{wq}, \code{J} & $(n_q,)$ & $(16,)$ o $(32,)$\\
\code{Pn}, \code{dP}, \code{W} & $(n_q,3)$ & $(16,3)$ o $(32,3)$\\
\code{X}, \code{dX}, normal de cuadratura & $(n_q,2)$ & $(16,2)$ o $(32,2)$\\
\code{U}, \code{T}, \code{Ts} & $(n_q,2,2)$ & Un tensor $2\times2$ por punto\\
\code{Gb}, \code{Hb} & $(3,2,2)$ & Un bloque por nodo local\\
\code{Hs}, \code{H}, \code{G} & $(2N,2N)$ & $(32,32)$\\
Normal nodal en \code{resolver} & $(N,2)$ & $(16,2)$\\
\code{t}, \code{G @ t} & $(2N,)$ & $(32,)$\\
\code{u}, tras \code{reshape} & $(N,2)$ & $(16,2)$\\
Resultado de \code{resolver} & $(N,)$ & $(16,)$ complejo\\
\bottomrule
\end{longtable}

\section{Constructor: guardar par\'ametros y crear la malla}
\label{sec:inicio}
\subsection{L\'ineas 1--3: clase, constructor y atributos}
\begin{lstlisting}
class BEMCilindro:
    def __init__(self, R, ne, m):
        self.R, self.ne, self.m = R, ne, m
\end{lstlisting}
\code{class} define el tipo de objeto. \code{__init__} es su
inicializador; no es una funci\'on que el usuario deba llamar
separadamente despu\'es de crear la instancia.

La asignaci\'on m\'ultiple equivale a:
\begin{lstlisting}
self.R = R
self.ne = ne
self.m = m
\end{lstlisting}
\code{m} es un diccionario con propiedades tales como
\code{m["nu"]}, \code{m["mu"]}, \code{m["lam"]},
\code{m["rho"]}, \code{m["cs"]} y \code{m["cp"]}.
La clase guarda una referencia al diccionario, no una copia independiente.

\begin{cuidado}
No se debe cambiar arbitrariamente la geometr\'ia o el material del
objeto despu\'es del prec\'alculo. Por ejemplo, modificar $\nu$
sin reconstruirlo deja $T^s$ y $H^s$ almacenados con el valor anterior.
Para otro problema es m\'as seguro crear otra instancia.
\end{cuidado}

\subsection{L\'inea 4: nodos y conectividad}
\begin{lstlisting}
self.nodes, self.conn = malla_circulo(R, ne)
\end{lstlisting}
La funci\'on devuelve dos arreglos. El primero contiene coordenadas;
el segundo indica qu\'e tres nodos globales forman cada elemento.
Para ocho elementos:
\begin{lstlisting}[language={}]
conn =
[[ 0,  1,  2],
 [ 2,  3,  4],
 [ 4,  5,  6],
 [ 6,  7,  8],
 [ 8,  9, 10],
 [10, 11, 12],
 [12, 13, 14],
 [14, 15,  0]]
\end{lstlisting}
Cada elemento tiene dos extremos compartidos y un nodo intermedio.
El \'ultimo elemento cierra el c\'irculo volviendo al nodo cero.

Los primeros tres nodos, en metros, son aproximadamente
\begin{equation}
 \code{nodes[:3]}=
 \begin{bmatrix}
 6&0\\
 5.54327720&2.29610059\\
 4.24264069&4.24264069
 \end{bmatrix}.
\end{equation}

\subsection{L\'inea 5: contar nodos}
\begin{lstlisting}
self.N = len(self.nodes)
\end{lstlisting}
\code{len} de un arreglo bidimensional devuelve el n\'umero de filas.
Por eso aqu\'i da 16, no 32. Hay 16 nodos, y cada uno tiene
dos componentes de desplazamiento: existen 32 inc\'ognitas escalares.

\subsection{L\'ineas 6--8: preparar almacenamiento}
La l\'inea 6 es un comentario; no ejecuta c\'alculos.
\begin{lstlisting}
self.quad = []
Hs = np.zeros((2*self.N, 2*self.N))
\end{lstlisting}
\code{self.quad} es una lista vac\'ia donde se guardar\'a informaci\'on
de cuadratura para reutilizarla. No es una matriz NumPy.

\code{Hs} es una matriz real de ceros, aqu\'i $32\times32$.
Es real porque el n\'ucleo de tracci\'on est\'atico usado depende de
la geometr\'ia y del Poisson real. Las matrices din\'amicas se
crear\'an como complejas m\'as adelante.

\begin{idea}
Esta l\'inea solo reserva espacio. Una matriz inicial de ceros no
significa que la matriz BEM final sea nula ni dispersa.
Sus bloques se ir\'an acumulando durante los bucles.
\end{idea}

\section{Constructor: recorrer pares nodo--elemento}
\subsection{L\'ineas 9--11: el significado de los bucles anidados}
\begin{lstlisting}
for ic, xi in enumerate(self.nodes):
    fila = []
    for cn in self.conn:
\end{lstlisting}
\code{enumerate} devuelve simult\'aneamente un contador y un valor:
\begin{center}
\begin{tabular}{lll}
\toprule
Iteraci\'on & \code{ic} & \code{xi}\\
\midrule
Primera & 0 & Coordenadas del nodo 0\\
Segunda & 1 & Coordenadas del nodo 1\\
$\cdots$ & $\cdots$ & $\cdots$\\
\'Ultima & 15 & Coordenadas del nodo 15\\
\bottomrule
\end{tabular}
\end{center}
Para cada uno de esos nodos, el bucle interior recorre los ocho elementos.
As\'i se consideran $16(8)=128$ pares.
\code{fila} se crea de nuevo para cada nodo y almacenar\'a los datos
de sus ocho pares con los elementos.

\begin{figure}[htbp]
\centering
\begin{tikzpicture}[>=Stealth,
 box/.style={draw,rounded corners,align=left,fill=blue!3,
 text width=11.6cm,minimum height=0.7cm},node distance=0.3cm]
 \node[box] (a){Fijar nodo de colocaci\'on: \code{ic = 0},
 $\vect\xi=(6,0)$};
 \node[box,below=of a] (b){Elemento 0: \code{cn = [0,1,2]}\\
 Contiene la fuente en el extremo izquierdo.};
 \node[box,below=of b] (c){Elemento 1: \code{cn = [2,3,4]}\\
 No contiene la fuente.};
 \node[box,below=of c] (d){$\cdots$};
 \node[box,below=of d] (e){Elemento 7: \code{cn = [14,15,0]}\\
 Contiene la fuente en el extremo derecho.};
 \node[box,below=of e] (f){Guardar esa fila; repetir ahora con
 \code{ic = 1}.};
 \foreach \s/\t in {a/b,b/c,c/d,d/e,e/f}
   \draw[->] (\s)--(\t);
\end{tikzpicture}
\caption{Una vuelta del bucle exterior. El nodo fuente permanece fijo
mientras se integran todos los elementos.}
\end{figure}

\subsection{L\'inea 12: buscar el nodo de colocaci\'on en el elemento}
\begin{lstlisting}
loc = np.where(cn == ic)[0]
\end{lstlisting}
Primero se compara cada entrada de \code{cn} con \code{ic}.
Con \code{ic=0} y \code{cn=[0,1,2]}:
\begin{lstlisting}[language={}]
cn == ic                 -> [True, False, False]
np.where(cn == ic)        -> (array([0]),)
np.where(cn == ic)[0]     -> array([0])
\end{lstlisting}
El primer \code{[0]} toma el arreglo de \'indices dentro de la
tupla que devuelve \code{np.where}. No selecciona todav\'ia
el primer \'indice contenido en ese arreglo.

Otros dos casos:
\begin{lstlisting}[language={}]
ic=0, cn=[2,3,4]    -> loc=array([])
ic=0, cn=[14,15,0]  -> loc=array([2])
\end{lstlisting}
Con la conectividad v\'alida del notebook, hay una coincidencia o ninguna.

\subsection{L\'inea 13: convertir posici\'on local en coordenada local}
\begin{lstlisting}
eta_s = [-1.0, 0.0, 1.0][loc[0]] if len(loc) else None
\end{lstlisting}
Se puede leer como:
\begin{lstlisting}
if len(loc):
    eta_s = [-1.0, 0.0, 1.0][loc[0]]
else:
    eta_s = None
\end{lstlisting}
Si \code{loc} est\'a vac\'io, no se eval\'ua \code{loc[0]},
evitando un error de \'indice. Las tres posibilidades locales son
\begin{equation}
 0\longmapsto-1,\qquad1\longmapsto0,\qquad2\longmapsto1.
\end{equation}
\code{None} significa que el elemento no contiene el nodo fuente;
no significa que la fuente est\'e en $\eta=0$.

\subsection{L\'inea 14: elegir la cuadratura adecuada}
\begin{lstlisting}
eta, wq = puntos_cuadratura(eta_s)
\end{lstlisting}
Se obtienen coordenadas locales y pesos.
La funci\'on auxiliar utiliza:
\begin{itemize}
 \item \code{None}: Gauss-Legendre ordinario, 16 puntos.
 \item $\eta_s=-1$ o $1$: mapeo cuadr\'atico hacia el extremo,
 16 puntos.
 \item $\eta_s=0$: dos subintervalos con 16 puntos cada uno,
 32 puntos en total.
\end{itemize}
No es una cuadratura adaptativa arbitraria: estas reglas est\'an
fijadas en la funci\'on auxiliar del notebook.

La fuente provoca singularidades al integrar su propio elemento.
El mapeo concentra puntos cerca de ella, pero los puntos de Gauss
no coinciden exactamente con la fuente.
El tratamiento de la singularidad fuerte requiere adem\'as
la construcci\'on especial de $H^s$ explicada m\'as adelante.

\section{Constructor: geometr\'ia y pesos de integraci\'on}
\subsection{L\'inea 15: evaluar la base y sus derivadas}
\begin{lstlisting}
Pn, dP = shape(eta), dshape(eta)
\end{lstlisting}
La fila $q$ de \code{Pn} contiene
\begin{equation}
 \left[P_1(\eta_q),P_2(\eta_q),P_3(\eta_q)\right],
\end{equation}
con
\begin{equation}
 P_1=\tfrac12\eta(\eta-1),\quad P_2=1-\eta^2,\quad
 P_3=\tfrac12\eta(\eta+1).
\end{equation}
\code{dP} contiene sus derivadas:
\begin{equation}
 \left[\eta_q-\tfrac12,\ -2\eta_q,\ \eta_q+\tfrac12\right].
\end{equation}
Ambos arreglos tienen forma $(n_q,3)$.
Estas derivadas son respecto a $\eta$, no respecto a $x$ o a $y$.

\subsection{L\'inea 16: multiplicar matrices para obtener la curva}
\begin{lstlisting}
X, dX = Pn @ self.nodes[cn], dP @ self.nodes[cn]
\end{lstlisting}
\code{self.nodes[cn]} selecciona las tres filas de coordenadas del
elemento; su forma es $(3,2)$.
El operador \code{@} es producto matricial:
\begin{equation}
 (n_q,3)\ @\ (3,2)\longrightarrow(n_q,2).
\end{equation}
En componentes:
\begin{align}
 X_{q,j}&=\sum_{a=0}^2P_a(\eta_q)x_{g(a),j},\\
 dX_{q,j}&=\sum_{a=0}^2P_a'(\eta_q)x_{g(a),j}.
\end{align}
Por tanto, \code{X[q]} es un punto sobre la curva interpolada y
\code{dX[q]} su tangente no normalizada.

\begin{cuidado}
\code{dX} no es la distancia desde la fuente hasta \code{X}.
Es la derivada geom\'etrica $\dd\vect x/\dd\eta$.
La distancia a la fuente la calculan los n\'ucleos a partir de
\code{X} y \code{xi}.
\end{cuidado}

\subsection{Ejemplo sencillo en el punto medio local}
Para el primer elemento, en $\eta=0$:
\begin{equation}
 P=[0,1,0],\qquad P'=[-1/2,0,1/2].
\end{equation}
Entonces
\begin{align}
 \vect X(0)&=\vect x_1=(5.54327720,\ 2.29610059)\ {\rm m},\\
 \vect x'(0)&=\frac{\vect x_2-\vect x_0}{2}
       =(-0.87867966,\ 2.12132034)\ {\rm m}.
\end{align}
Este es un punto did\'actico para comprobar la interpolaci\'on;
no se afirma que $\eta=0$ sea uno de los puntos de Gauss utilizados.

\subsection{L\'inea 17: jacobiano por cada fila}
\begin{lstlisting}
J = np.linalg.norm(dX, axis=-1)
\end{lstlisting}
En un arreglo $(n_q,2)$, el \'ultimo eje es el de las dos componentes.
Por eso:
\begin{equation}
 J_q=\sqrt{dX_{q,0}^2+dX_{q,1}^2}.
\end{equation}
La forma de salida es $(n_q,)$, no un solo n\'umero.
F\'isicamente:
\begin{equation}
 \dd\Gamma=J(\eta)\dd\eta.
\end{equation}
El jacobiano convierte un intervalo local adimensional en una longitud
del elemento f\'isico.

\subsection{L\'inea 18: normal exterior y \textit{broadcasting}}
\begin{lstlisting}
nrm = np.stack([dX[:, 1], -dX[:, 0]], -1)/J[:, None]
\end{lstlisting}
Se realizan tres operaciones:
\begin{enumerate}
 \item \code{dX[:,1]} toma todas las derivadas $y'$;
 \code{-dX[:,0]} toma todas las derivadas $-x'$.
 \item \code{np.stack(...,-1)} las coloca como dos columnas:
 $(y',-x')$.
 \item Se divide cada fila por su longitud $J_q$.
\end{enumerate}
Para un contorno antihorario:
\begin{equation}
 \vect n_q=\frac{(y_q',-x_q')}{J_q}.
\end{equation}
La rotaci\'on hacia la derecha da la normal exterior.

\code{J[:,None]} transforma la forma $(n_q,)$ en $(n_q,1)$:
\begin{equation}
 \begin{bmatrix}J_0\\J_1\\\vdots\end{bmatrix}.
\end{equation}
NumPy extiende esa columna sobre las dos componentes para efectuar la
divisi\'on. Eso es \emph{broadcasting}; no hay una suma oculta.
Sin \code{None}, las formas no describir\'ian la divisi\'on
por filas pretendida.

\begin{figure}[htbp]
\centering
\begin{tikzpicture}[>=Stealth,scale=0.95]
 \draw[thick,blue] (0,0.5)..controls(1.7,-0.8)and(3,-0.8)..(4.5,0.4);
 \fill (2.2,-0.48)circle(2pt);
 \draw[->,green!40!black,thick] (2.2,-0.48)--(3.6,-0.48)
 node[right]{$(x',y')/J$};
 \draw[->,red!65!black,thick] (2.2,-0.48)--(2.2,-1.7)
 node[below]{$(y',-x')/J$};
 \node[align=left] at (7.4,-0.4){Tangente hacia la derecha\\
 Interior del cuerpo arriba\\Normal exterior hacia abajo};
\end{tikzpicture}
\caption{Orientaci\'on en un tramo inferior de una frontera antihoraria.
Invertir el recorrido sin cambiar la f\'ormula invierte la normal.}
\end{figure}

\subsection{L\'inea 19: combinar los factores que se repetir\'an}
\begin{lstlisting}
W = Pn*(J*wq)[:, None]
\end{lstlisting}
Aqu\'i \code{*} es producto entrada por entrada, no producto matricial.
Primero se forma $J_qw_q$. Luego \code{[:,None]} lo convierte en
columna y se multiplica por las tres funciones de forma:
\begin{equation}
 \boxed{W_{qa}=P_a(\eta_q)J_qw_q.}
 \label{eq:W}
\end{equation}
As\'i la integral de un n\'ucleo $K$ queda aproximadamente
\begin{equation}
 \int_{-1}^1K(\eta)P_a(\eta)J(\eta)\dd\eta
       \simeq\sum_qW_{qa}K(\eta_q).
\end{equation}
Los pesos \code{wq} ya incluyen, cuando corresponde, el cambio de
variable de la cuadratura singular. No hay que multiplicar
otra vez por ese cambio de variable.

Aunque los pesos de Gauss sean positivos, alguna entrada de \code{W}
puede ser negativa porque una funci\'on de forma cuadr\'atica
puede ser negativa entre nodos.

\subsection{Una fila de cuadratura real, comprobada}
Con \code{ic=0}, primer elemento y primer punto transformado:
\begin{align}
 \eta_0&\simeq-0.999943829910,\\
 w_0&\simeq2.87790682442\times10^{-4},\\
 P(\eta_0)&\simeq
 [0.9999157464,\ 1.123370240\times10^{-4},\
   -2.808346722\times10^{-5}],\\
 \vect X_0&\simeq
 (5.9999980459,\ 1.387890468\times10^{-4})\ {\rm m},\\
 \vect x'_0&\simeq(-0.0348133561,\ 2.4708612100)\ {\rm m},\\
 J_0&\simeq2.4711064503\ {\rm m},\\
 \vect n_0&\simeq(0.9999007569,\ 0.0140881653),\\
 W_{0,:}&\simeq
 [7.111014938\times10^{-4},\
  7.988975656\times10^{-8},\
 -1.997187820\times10^{-8}]\ {\rm m}.
\end{align}
Se observa directamente que el punto est\'a muy cerca de la fuente
$(6,0)$, sin coincidir con ella.
La normal es la de la curva cuadr\'atica aproximada, no necesariamente
la normal radial exacta del c\'irculo.

\section{Constructor: integrar Kelvin y completar la diagonal}
\label{sec:estatica}
\subsection{L\'inea 20: evaluar el n\'ucleo, no la matriz global}
\begin{lstlisting}
Ts = kernel_Ts(X, xi, nrm, m["nu"])
\end{lstlisting}
Por cada punto \code{X[q]} se eval\'ua un bloque:
\begin{equation}
 \code{Ts[q,i,j]}=T^s_{ij}(\vect X_q,\vect\xi,\vect n_q).
\end{equation}
La forma de \code{Ts} es $(n_q,2,2)$.
\code{m["nu"]} suministra el Poisson.
El n\'ucleo est\'atico de tracci\'on utilizado no depende de la
frecuencia; de ah\'i que pueda guardarse.

\subsection{L\'inea 21: pasar de nodo local a nodo global}
\begin{lstlisting}
for a, g in enumerate(cn):
\end{lstlisting}
Con \code{cn=[2,3,4]} se obtiene:
\begin{center}
\begin{tabular}{ccc}
\toprule
\code{a} & \code{g} & Lectura\\
\midrule
0 & 2 & Primer nodo local, nodo global 2\\
1 & 3 & Segundo nodo local, nodo global 3\\
2 & 4 & Tercer nodo local, nodo global 4\\
\bottomrule
\end{tabular}
\end{center}
\code{a} selecciona una funci\'on de forma en \code{W}.
\code{g} selecciona un bloque de columnas en la matriz global.
Intercambiarlos produce un ensamblaje incorrecto.

\subsection{L\'inea 22: posponer el bloque diagonal}
\begin{lstlisting}
if g != ic:
\end{lstlisting}
El bloque diagonal global corresponde a que el nodo de columna sea
el mismo nodo de colocaci\'on.
Se omite \emph{solo por ahora} porque contiene la contribuci\'on
singular fuerte y el t\'ermino libre.
No se est\'a suponiendo que ese bloque sea cero.

Los bloques no diagonales que proceden de elementos con fuente
tienen una funci\'on de forma que se anula en ella, lo cual compensa
la singularidad principal del integrando.

\subsection{L\'inea 23: entender primero la suma de cuadratura}
\begin{lstlisting}
np.einsum("q,qij->ij", W[:, a], Ts)
\end{lstlisting}
Se lee:
\begin{equation}
 B_{ij}=\sum_qW_{qa}T^s_{qij}.
 \label{eq:einsum-estatico}
\end{equation}
\code{W[:,a]} tiene forma $(n_q,)$ y \code{Ts} forma $(n_q,2,2)$.
El corte \code{W[:,a]} toma todas las filas y solo la columna del nodo
local \code{a}.
La letra \code{q} desaparece despu\'es de la flecha:
ese es el \'indice que se suma.
Las letras \code{i,j} permanecen y producen un bloque $(2,2)$.

Un equivalente explicativo es:
\begin{lstlisting}
bloque = np.zeros((2, 2))
for q in range(len(W)):
    bloque += W[q, a]*Ts[q]
\end{lstlisting}
No hay un determinante ni una inversi\'on escondidos en \code{einsum}.
Es una contracci\'on de \'indices.

\subsection{L\'inea 23: entender despu\'es los cortes de la matriz}
\begin{lstlisting}
Hs[2*ic:2*ic+2, 2*g:2*g+2] += bloque
\end{lstlisting}
En Python, el extremo superior de un corte no se incluye.
Por ejemplo, con \code{ic=3} y \code{g=5}:
\begin{lstlisting}
Hs[6:8, 10:12]
\end{lstlisting}
selecciona filas 6 y 7, columnas 10 y 11:
un bloque de dos por dos.
La regla proviene del orden de componentes
\begin{equation}
 (u_{0x},u_{0y},u_{1x},u_{1y},\ldots).
\end{equation}
El operador \code{+=} suma una contribuci\'on al valor ya almacenado.
No debe sustituirse por \code{=}, porque nodos compartidos reciben
contribuciones de m\'as de un elemento.

\begin{figure}[htbp]
\centering
\begin{tikzpicture}[scale=0.78]
 \draw[step=0.5,gray!45] (0,0)grid(6,6);
 \foreach \v in {0,1,...,6}{
   \draw[thick] (\v,0)--(\v,6);
   \draw[thick] (0,\v)--(6,\v);}
 \fill[blue!30,opacity=0.7] (4,2)rectangle(5,3);
 \draw[very thick,blue!70!black] (4,2)rectangle(5,3);
 \node[above] at (4.5,6.15){Nodo global de columna $g$};
 \draw[->,thick] (4.5,6.05)--(4.5,3.15);
 \node[align=right,left] at (-0.25,2.5){Nodo de\\colocaci\'on $c$};
 \node at (7.1,2.5){Bloque $2\times2$};
\end{tikzpicture}
\caption{Cada bloque conecta dos ecuaciones de un nodo con dos
componentes de otro. El esquema representa bloques, no un tama\~no
concreto de malla.}
\end{figure}

\subsection{L\'ineas 24--25: almacenar el prec\'alculo}
\begin{lstlisting}
fila.append((cn, X, nrm, W, Ts))
self.quad.append(fila)
\end{lstlisting}
La primera instrucci\'on est\'a dentro del bucle de elementos:
agrega una tupla de cinco objetos para un par concreto.
La segunda est\'a fuera de ese bucle y dentro del bucle de nodos:
agrega la fila completa.

La estructura puede imaginarse as\'i:
\begin{lstlisting}[language={}]
self.quad
  [0]  -> datos de todos los elementos para la fuente 0
      [0] -> (cn, X, nrm, W, Ts) del elemento 0
      [1] -> (cn, X, nrm, W, Ts) del elemento 1
      ...
  [1]  -> datos de todos los elementos para la fuente 1
      ...
\end{lstlisting}
Con ocho elementos:
\begin{equation}
 \code{len(self.quad)}=16,\qquad
 \code{len(self.quad[ic])}=8.
\end{equation}
Hay 128 tuplas, no 128 nodos.

\begin{cuidado}
\code{fila} es una fila de la \emph{lista de datos de cuadratura}.
No es una fila num\'erica de $H$.
Adem\'as, los datos se guardan por par nodo--elemento, no solamente
por elemento: la cuadratura depende de d\'onde est\'a la fuente.
\end{cuidado}

\subsection{Cu\'antos puntos se guardan en total}
En esta malla, cada nodo extremo pertenece a dos elementos:
ambos usan 16 puntos; los otros seis tambi\'en usan 16.
Esa fila almacena $8(16)=128$ puntos.

Cada nodo intermedio pertenece a un elemento que usa 32 puntos,
mientras que los otros siete usan 16:
$32+7(16)=144$.
Como hay ocho nodos de cada tipo:
\begin{equation}
 8(128)+8(144)=2176
\end{equation}
evaluaciones geom\'etricas almacenadas, contando todos los pares.
No son 2176 puntos f\'isicos necesariamente distintos.

\subsection{L\'ineas 26--27: construir los bloques diagonales}
\begin{lstlisting}
for ic in range(self.N):
    Hs[2*ic:2*ic+2, 2*ic:2*ic+2] = -sum(
        Hs[2*ic:2*ic+2, 2*g:2*g+2]
        for g in range(self.N) if g != ic
    )
\end{lstlisting}
El salto de l\'inea anterior es solo para facilitar la lectura.
La expresi\'on usa \code{sum} de Python sobre un generador de
matrices $2\times2$.
Para cada fila de bloques:
\begin{equation}
 \boxed{H^s_{cc}=-\sum_{g\ne c}H^s_{cg}.}
 \label{eq:diagonal}
\end{equation}
\textbf{Justificaci\'on:} una traslaci\'on est\'atica constante no
deforma el cuerpo y no genera tracciones. Como las funciones de forma
suman uno, el BEM est\'atico debe satisfacer
\begin{equation}
 \sum_gH^s_{cg}=0.
\end{equation}
Se despeja el bloque diagonal y se obtiene \eqref{eq:diagonal}.
Esta operaci\'on incorpora el t\'ermino libre y el valor principal
est\'atico sin calcularlos por separado.

\begin{cuidado}[Dos errores muy frecuentes]
La suma se hace por \emph{bloques} $2\times2$, no sumando todas las
entradas escalares de una fila.
Tampoco debe agregarse un $I/2$ adicional: ya est\'a incorporado
en esta construcci\'on de la diagonal.
\end{cuidado}

En la ejecuci\'on del ejemplo:
\begin{equation}
 H^s_{00}\simeq
 \begin{bmatrix}
 0.492760151&0\\
 0&0.461415868
 \end{bmatrix}.
 \label{eq:diag-ejemplo}
\end{equation}
Las entradas fuera de esa diagonal son del orden de $10^{-16}$.
Este bloque no es exactamente $I/2$: contiene m\'as que el t\'ermino
libre aislado.

\subsection{L\'inea 28: conservar la matriz para otras frecuencias}
\begin{lstlisting}
self.Hs = Hs
\end{lstlisting}
El nombre local \code{Hs} se guarda como atributo.
Al terminar el constructor, los datos importantes quedan en el objeto:
\begin{equation}
 \code{R},\ \code{ne},\ \code{m},\ \code{nodes},\
 \code{conn},\ \code{N},\ \code{quad},\ \code{Hs}.
\end{equation}
La l\'inea 29 est\'a en blanco: separa visualmente m\'etodos,
sin efectuar ninguna operaci\'on.

\section{\texttt{matrices(w)}: ensamblar una frecuencia}
\label{sec:matrices}
\subsection{L\'ineas 30--31: entrada y abreviatura local}
\begin{lstlisting}
def matrices(self, w):
    N = self.N
\end{lstlisting}
\code{w} es una frecuencia angular escalar. La clase no convierte
autom\'aticamente de Hz a rad/s.
\code{N} es una abreviatura local para no escribir repetidamente
\code{self.N}; no vuelve a contar los nodos ni cambia la malla.

\subsection{L\'inea 32: empezar \texorpdfstring{$G$}{G} desde cero}
\begin{lstlisting}
G = np.zeros((2*N, 2*N), complex)
\end{lstlisting}
Se reserva una matriz compleja de $32\times32$ en el ejemplo.
Se integrar\'a $U$ para construirla. Se necesita tipo complejo porque
los n\'ucleos arm\'onicos contienen amplitud y fase.

\subsection{L\'inea 33: partir de una copia compleja de \texorpdfstring{$H^s$}{Hs}}
\begin{lstlisting}
H = self.Hs.astype(complex)
\end{lstlisting}
La matriz \code{self.Hs} es real. La conversi\'on crea un arreglo
complejo separado. Por eso las sumas posteriores sobre \code{H}
no cambian la matriz est\'atica guardada.

El esquema es
\begin{equation}
 H(\omega)=H^s+\Delta H(\omega),\qquad
 \Delta H^e_{ca}\simeq\sum_qW_{qa}
       [T(\vect X_q,\vect\xi_c,\omega)-T^s(\vect X_q,\vect\xi_c)].
 \label{eq:correccion}
\end{equation}
Se comprob\'o que \code{matrices(40.0)} conserva \code{self.Hs}.

\begin{cuidado}
Si se escribiera \code{H = self.Hs}, ambos nombres se referir\'ian
al mismo arreglo. No ser\'ia una copia.
Adem\'as de los problemas de tipo al sumar complejos, se podr\'ia
corromper el prec\'alculo entre frecuencias.
\end{cuidado}

\subsection{L\'ineas 34--35: recuperar datos ya preparados}
\begin{lstlisting}
for ic, xi in enumerate(self.nodes):
    for cn, X, nrm, W, Ts in self.quad[ic]:
\end{lstlisting}
Se recorre nuevamente cada nodo de colocaci\'on.
Dentro se leen las tuplas guardadas para esa fuente.
El desempaquetado asigna sus cinco objetos a los cinco nombres.

No se vuelven a calcular \code{shape}, \code{dshape},
\code{J} ni la cuadratura. Esa es la utilidad del constructor:
evitar repetir geometr\'ia para cada frecuencia.
Los \'indices de los puntos de cuadratura est\'an impl\'icitos
en las filas de los arreglos.

\subsection{L\'inea 36: evaluar los n\'ucleos din\'amicos}
\begin{lstlisting}
U, T = kernels(X, xi, nrm, w, self.m)
\end{lstlisting}
Se eval\'uan:
\begin{equation}
 U_{qij}=U_{ij}(\vect X_q,\vect\xi,\omega),\qquad
 T_{qij}=T_{ij}(\vect X_q,\vect\xi,\vect n_q,\omega).
\end{equation}
Ambos arreglos tienen forma $(n_q,2,2)$.
La frecuencia y las propiedades complejas del material entran en
esta operaci\'on.
$U$ no depende de la normal, pero la funci\'on recibe la normal
porque tambi\'en calcula $T$.

\subsection{L\'inea 37: integrar los tres bloques de \texorpdfstring{$G$}{G}}
\begin{lstlisting}
Gb = np.einsum("qa,qij->aij", W, U)
\end{lstlisting}
La expresi\'on significa
\begin{equation}
 \boxed{(G_b)_{aij}=\sum_qW_{qa}U_{qij}.}
 \label{eq:Gb}
\end{equation}
El \'indice \code{q} desaparece; \code{a,i,j} permanecen.
Por tanto, la salida tiene forma $(3,2,2)$.
No se suman los tres nodos locales: se conservan tres bloques distintos
porque cada uno debe ir a una columna global diferente.

Un equivalente para comprenderlo:
\begin{lstlisting}
Gb = np.zeros((3, 2, 2), complex)
for a in range(3):
    for q in range(len(W)):
        Gb[a] += W[q, a]*U[q]
\end{lstlisting}

\subsection{L\'inea 38: integrar la correcci\'on din\'amica}
\begin{lstlisting}
Hb = np.einsum("qa,qij->aij", W, T - Ts)
\end{lstlisting}
La operaci\'on es la misma que en \eqref{eq:Gb}, con otro n\'ucleo:
\begin{equation}
 (H_b)_{aij}=\sum_qW_{qa}(T_{qij}-T^s_{qij}).
\end{equation}
Se resta $T^s$ porque la singularidad principal de $T$ coincide
con la de Kelvin. La parte est\'atica ya est\'a en \code{self.Hs}.

\begin{idea}[Por qu\'e no se integra simplemente $T$ y se suma a $H^s$]
Eso contar\'ia dos veces la parte est\'atica.
La separaci\'on correcta es
$T=T^s+(T-T^s)$: una parte en la matriz precalculada y la otra en
la correcci\'on din\'amica.
\end{idea}

La diferencia es regularizada en sentido anal\'itico. Esto no significa
que restar los n\'ucleos en doble precisi\'on sea siempre inocuo:
a distancias extremadamente peque\~nas puede haber cancelaci\'on.
La gu\'ia general contiene el estudio num\'erico de esa limitaci\'on.

\subsection{L\'ineas 39--41: sumar los bloques en su posici\'on global}
\begin{lstlisting}
for a, g in enumerate(cn):
    G[2*ic:2*ic+2, 2*g:2*g+2] += Gb[a]
    H[2*ic:2*ic+2, 2*g:2*g+2] += Hb[a]
\end{lstlisting}
Son las mismas reglas local--global del constructor.
Hay una diferencia importante: aqu\'i no aparece
\code{if g != ic}. Se incluyen tambi\'en los bloques diagonales,
porque se est\'a integrando la diferencia regularizada $T-T^s$.

\begin{center}
\begin{tabular}{lll}
\toprule
Momento & N\'ucleo & Tratamiento del bloque diagonal\\
\midrule
Constructor & $T^s$ & Se omite y reconstruye por traslaci\'on\\
\code{matrices} & $T-T^s$ & Se integra y suma normalmente\\
\code{matrices} & $U$ & Se integra con la cuadratura preparada\\
\bottomrule
\end{tabular}
\end{center}

\subsection{Ejemplo real de dos elementos que comparten un nodo}
Fijemos \code{ic=0} y el nodo global de columna \code{g=2}.
Ese nodo pertenece a:
\begin{itemize}
 \item Elemento 0, \code{cn=[0,1,2]}, con \code{a=2}.
 \item Elemento 1, \code{cn=[2,3,4]}, con \code{a=0}.
\end{itemize}
Por tanto:
\begin{equation}
 G_{0,2}=G^{e=0}_{b,\,a=2}+G^{e=1}_{b,\,a=0}.
\end{equation}
En el ejemplo, la entrada superior izquierda de esos bloques es
\begin{align}
 (G^{e=0}_{b,2})_{00}
 &\simeq-8.384184254\times10^{-8}
         -\ii\,2.994329428\times10^{-8},\\
 (G^{e=1}_{b,0})_{00}
 &\simeq-6.683467359\times10^{-8}
         -\ii\,2.253818621\times10^{-8}.
\end{align}
La entrada ensamblada es la suma:
\begin{equation}
 G_{0,2;\,00}\simeq
 -1.506765161\times10^{-7}
 -\ii\,5.248148049\times10^{-8}\ {\rm m/Pa}.
\end{equation}
En \'indices escalares de NumPy, el bloque completo es
\code{G[0:2,4:6]}; su entrada superior izquierda es \code{G[0,4]}.
Se comprob\'o num\'ericamente el bloque completo, no solo esa entrada.

\subsection{L\'inea 42: devolver dos matrices, no desplazamientos}
\begin{lstlisting}
return H, G
\end{lstlisting}
Python devuelve una tupla de dos arreglos.
Todav\'ia no se aplic\'o la presi\'on en este m\'etodo.
Cambiar $p$ no cambia las matrices si se mantienen geometr\'ia,
material y frecuencia: solo cambia el segundo miembro.
La l\'inea 43 es un separador en blanco.

\section{\texttt{resolver(w,p)}: carga y sistema lineal}
\label{sec:resolver}
\subsection{L\'inea 44: frecuencia y presi\'on predeterminada}
\begin{lstlisting}
def resolver(self, w, p=P):
\end{lstlisting}
El usuario debe indicar \code{w}. Puede omitir \code{p};
en ese caso se usa el valor de \code{P} existente cuando se defini\'o
el m\'etodo.

\begin{lstlisting}
ur = bem.resolver(40.0)
ur = bem.resolver(40.0, p=100.0)
\end{lstlisting}
Para la definici\'on original, ambas llamadas prescriben 100 Pa.
En Python, los argumentos predeterminados se fijan al ejecutar la
definici\'on, no se releen de la variable global en cada llamada.
Para otra carga, pasar \code{p} expl\'icitamente evita ambig\"uedades.

\subsection{L\'inea 45: construir las matrices de esa frecuencia}
\begin{lstlisting}
H, G = self.matrices(w)
\end{lstlisting}
Esta llamada ejecuta todo el ensamblaje explicado en la
secci\'on~\ref{sec:matrices}. El constructor no se ejecuta otra vez.
Se reutilizan sus datos precalculados, pero no se almacena una cach\'e
de las matrices din\'amicas entre llamadas.

\subsection{L\'inea 46: normal radial en cada nodo}
\begin{lstlisting}
nrm = self.nodes/self.R
\end{lstlisting}
Cada nodo est\'a sobre el c\'irculo de radio $R$:
\begin{equation}
 \vect n_g=\frac{\vect x_g}{R}
          =(\cos\theta_g,\sin\theta_g).
\end{equation}
El resultado tiene forma $(N,2)$.

\begin{cuidado}[La variable \texttt{nrm} tiene dos usos diferentes]
En el constructor y en \code{matrices}, \code{nrm} contiene normales
en puntos de \emph{cuadratura}, obtenidas de la curva interpolada.
En \code{resolver}, contiene normales \emph{nodales} del c\'irculo exacto.
Los nombres coinciden, pero los arreglos, su tama\~no y su papel son distintos.
\end{cuidado}

La normal nodal es apropiada para este c\'irculo centrado en el origen.
Dividir coordenadas por \code{R} no es una receta general para obtener
normales de cualquier geometr\'ia.
Entre nodos, la presi\'on interpolada a partir de estos valores no es
exactamente $-p$ por la normal geom\'etrica de la curva aproximada.

\subsection{L\'inea 47: de presi\'on escalar a vector global de tracciones}
\begin{lstlisting}
t = (-p*nrm).ravel()
\end{lstlisting}
Primero, \code{-p*nrm} crea un arreglo $(N,2)$:
cada fila contiene $(t_x,t_y)=-p(n_x,n_y)$.
\code{ravel()} lo aplana por filas en el orden predeterminado:
\begin{equation}
 \vect t=(t_{0x},t_{0y},t_{1x},t_{1y},\ldots)^{\mathsf T}.
\end{equation}
No se est\'an sumando ni promediando componentes.

Los primeros valores del ejemplo son
\begin{equation}
 \vect t_{0:6}\simeq
 [-100,\ 0,\ -92.38795325,\ -38.26834324,\
    -70.71067812,\ -70.71067812]^{\mathsf T}\ {\rm Pa}.
\end{equation}
Que la presi\'on sea uniforme no significa que todas las componentes
cartesianas de tracci\'on sean iguales: la normal cambia alrededor del borde.

\subsection{L\'inea 48, primera parte: formar el segundo miembro}
\begin{lstlisting}
G @ t
\end{lstlisting}
Las formas son
\begin{equation}
 (2N,2N)\ @\ (2N,)\longrightarrow(2N,).
\end{equation}
Se obtiene $\vect b=G\vect t$.
La tracci\'on tiene unidades Pa y $G$ unidades m/Pa, por lo que
$\vect b$ tiene unidades m. $H$ es adimensional.

Para el ejemplo, las primeras dos entradas son aproximadamente
\begin{equation}
 \vect b_{0:2}=
 \begin{bmatrix}
 -6.98841584\times10^{-5}+\ii\,9.22332748\times10^{-5}\\
 0
 \end{bmatrix}\ {\rm m}.
\end{equation}
La segunda entrada es del orden de $10^{-18}$ m, no exactamente cero
en aritm\'etica de punto flotante.
$\vect b$ no es a\'un el desplazamiento: debe resolverse el sistema.

\subsection{L\'inea 48, segunda parte: resolver, no invertir expl\'icitamente}
\begin{lstlisting}
np.linalg.solve(H, G @ t)
\end{lstlisting}
Busca el vector $\vect u$ que satisface $H\vect u=\vect b$.
No calcula una matriz inversa expl\'icita.
La salida tiene $2N$ componentes complejas, en el mismo orden
de nodos y direcciones que se us\'o al ensamblar.

La clase no agrega restricciones ni valores prescritos de desplazamiento:
en este problema todas las tracciones son conocidas.
Si el sistema fuera singular, el solucionador podr\'ia generar una
excepci\'on; el c\'odigo no la intercepta ni inventa una soluci\'on alternativa.

\subsection{L\'inea 48, tercera parte: restaurar filas de dos componentes}
\begin{lstlisting}
.reshape(-1, 2)
\end{lstlisting}
\code{-1} indica a NumPy que deduzca cu\'antas filas hacen falta.
Como hay $2N$ valores y dos columnas, resultan $N$ filas:
\begin{equation}
 [u_{0x},u_{0y},u_{1x},u_{1y},\ldots]
 \longrightarrow
 \begin{bmatrix}
 u_{0x}&u_{0y}\\
 u_{1x}&u_{1y}\\
 \vdots&\vdots
 \end{bmatrix}.
\end{equation}
No se resuelve otro sistema ni se cambian valores: solo se reorganiza
su forma.

Para el nodo cero del ejemplo:
\begin{align}
 u_{0x}&\simeq-2.395965868\times10^{-4}
                 +\ii\,4.085472003\times10^{-5}\ {\rm m},\\
 u_{0y}&\simeq2.883762378\times10^{-10}
                 +\ii\,2.193549058\times10^{-11}\ {\rm m}.
\end{align}
La componente $y$ es peque\~na frente a la radial, como se espera
cerca del eje $x$, pero no exactamente cero en el c\'alculo discreto.

\subsection{L\'inea 49: proyectar radialmente}
\begin{lstlisting}
return (u*nrm).sum(-1)
\end{lstlisting}
Primero se multiplican entradas correspondientes:
\begin{equation}
 \code{u*nrm}=
 \begin{bmatrix}
 u_{0x}n_{0x}&u_{0y}n_{0y}\\
 u_{1x}n_{1x}&u_{1y}n_{1y}\\
 \vdots&\vdots
 \end{bmatrix}.
\end{equation}
Despu\'es, \code{sum(-1)} suma el \'ultimo eje, las dos columnas:
\begin{equation}
 \boxed{u_{r,g}=u_{x,g}n_{x,g}+u_{y,g}n_{y,g}.}
 \label{eq:proyeccion}
\end{equation}
La salida tiene forma $(N,)$ y sigue siendo compleja.
No hay conjugaci\'on, valor absoluto ni promedio en esa l\'inea.

\begin{figure}[htbp]
\centering
\begin{tikzpicture}[>=Stealth,scale=1.05]
 \draw[->,gray] (0,0)--(4,0)node[right]{$x$};
 \draw[->,gray] (0,0)--(0,3)node[above]{$y$};
 \draw[->,thick,blue] (0,0)--(3,2)node[right]{$\vect n$};
 \draw[->,thick,red!65!black] (0,0)--(0.9,2.4)node[above]{$\vect u$};
 \draw[dashed] (0.9,2.4)--(1.7308,1.1538);
 \draw[->,very thick,green!40!black] (0,0)--(1.7308,1.1538)
 node[below right]{$u_r\vect n$};
 \node[align=left] at (6.6,1.5){Proyecci\'on geom\'etrica\\
 $\vect u\cdot\vect n=u_r$\\
 Aplicada tambi\'en a amplitudes complejas};
\end{tikzpicture}
\caption{Interpretaci\'on de la proyecci\'on. El dibujo ilustra
vectores reales; en el c\'odigo las dos componentes de desplazamiento
son complejas y la normal es real.}
\end{figure}

\section{Qu\'e recibe finalmente el usuario}
\subsection{Valores radiales, amplitudes y fases}
Los primeros resultados calculados son aproximadamente
\begin{equation}
 \vect u_r\simeq
 \begin{bmatrix}
 -2.395965868\times10^{-4}+\ii\,4.085472003\times10^{-5}\\
 -2.396737370\times10^{-4}+\ii\,4.087055508\times10^{-5}\\
 -2.395968718\times10^{-4}+\ii\,4.085489956\times10^{-5}\\
 \vdots
 \end{bmatrix}{\rm m}.
\end{equation}
Para obtener cantidades diferentes:
\begin{lstlisting}
ur = bem.resolver(40.0)
amplitud_por_nodo = np.abs(ur)
fase_por_nodo = np.angle(ur)
amplitud_media = np.abs(ur).mean()
instantanea_t0 = ur.real
\end{lstlisting}
\code{np.angle} devuelve radianes salvo que se indique
\code{deg=True}. La parte real es el desplazamiento f\'isico
en $t=0$ para la convenci\'on temporal del notebook.

\begin{cuidado}
No es lo mismo \code{np.abs(ur).mean()} que
\code{np.abs(ur.mean())}.
La primera expresi\'on promedia amplitudes; la segunda toma
la amplitud de un promedio complejo, que puede contener cancelaci\'on
de fases.
\end{cuidado}

\subsection{Comparaci\'on num\'erica con la referencia}
En el ejemplo:
\begin{align}
 \operatorname{media}|u_r^{\rm BEM}|
       &\simeq2.4309427051\times10^{-4}\ {\rm m},\\
 u_r^{\rm an}
       &\simeq-2.3986692491\times10^{-4}
                   +\ii\,4.0928186578\times10^{-5}\ {\rm m},\\
 |u_r^{\rm an}|&\simeq2.4333363541\times10^{-4}\ {\rm m}.
\end{align}
El error relativo de amplitud promedio es
\begin{equation}
 e_A=
 \frac{|\operatorname{media}|u_r^{\rm BEM}|-|u_r^{\rm an}||}
 {|u_r^{\rm an}|}
 \simeq9.83690\times10^{-4}
 =0.0983690\,\%.
\end{equation}
No es una medida completa del error de fase ni del vector cartesiano.
Se utiliza aqu\'i para comprobar que la lectura del c\'odigo coincide
con la salida que muestra el notebook.

\section{Resumen de todas las l\'ineas}
\begin{longtable}{p{0.13\textwidth}p{0.77\textwidth}}
\toprule
L\'ineas & Operaci\'on\\
\midrule
\endhead
1 & Definir la clase.\\
2 & Definir el inicializador con radio, elementos y material.\\
3 & Guardar par\'ametros como atributos.\\
4 & Crear nodos y conectividad.\\
5 & Guardar el n\'umero de nodos.\\
6 & Comentario que anuncia el prec\'alculo.\\
7 & Crear la lista global de datos de cuadratura.\\
8 & Reservar la matriz est\'atica real.\\
9 & Recorrer nodos de colocaci\'on y sus coordenadas.\\
10 & Crear una lista para los elementos vistos desde esa fuente.\\
11 & Recorrer las conectividades de los elementos.\\
12 & Buscar si la fuente es uno de sus nodos.\\
13 & Obtener su coordenada local o \code{None}.\\
14 & Generar puntos y pesos adecuados al par.\\
15 & Evaluar funciones de forma y derivadas.\\
16 & Obtener curva y tangente en puntos de cuadratura.\\
17 & Calcular el jacobiano en cada punto.\\
18 & Calcular normales exteriores unitarias.\\
19 & Combinar funciones de forma, jacobiano y pesos.\\
20 & Evaluar el n\'ucleo est\'atico de tracci\'on.\\
21 & Recorrer correspondencias local--global.\\
22 & Posponer el bloque diagonal est\'atico.\\
23 & Integrar y sumar un bloque no diagonal de $H^s$.\\
24 & Guardar la tupla del par nodo--elemento.\\
25 & Guardar la lista completa para esa fuente.\\
26 & Recorrer todas las diagonales globales.\\
27 & Reconstruir cada diagonal mediante traslaci\'on r\'igida.\\
28 & Guardar $H^s$ en el objeto.\\
29 & Separador visual, sin ejecuci\'on.\\
30 & Definir el m\'etodo para matrices din\'amicas.\\
31 & Crear una abreviatura local para el n\'umero de nodos.\\
32 & Reservar $G$ compleja y nula.\\
33 & Copiar $H^s$ a una matriz compleja.\\
34 & Recorrer nodos de colocaci\'on.\\
35 & Recuperar los datos guardados de sus elementos.\\
36 & Evaluar $U$ y $T$ para esa frecuencia.\\
37 & Integrar los tres bloques locales de $G$.\\
38 & Integrar los tres bloques de la correcci\'on $T-T^s$.\\
39 & Recorrer nuevamente los tres nodos locales.\\
40 & Acumular bloques de $G$ en posiciones globales.\\
41 & Acumular correcciones en $H$, incluidas las diagonales.\\
42 & Entregar ambas matrices.\\
43 & Separador visual, sin ejecuci\'on.\\
44 & Definir el solucionador con presi\'on predeterminada.\\
45 & Construir las matrices de la frecuencia solicitada.\\
46 & Obtener normales radiales nodales.\\
47 & Construir y aplanar las tracciones conocidas.\\
48 & Formar el segundo miembro, resolver y reorganizar componentes.\\
49 & Proyectar radialmente y devolver un valor complejo por nodo.\\
\bottomrule
\end{longtable}

\section{Comprobaciones realizadas y l\'imites de interpretaci\'on}
\label{sec:verificacion}
\subsection{Qu\'e se comprob\'o realmente}
Se ejecutaron las definiciones originales del notebook para la malla
de ocho elementos y 40 rad/s. Se verific\'o:
\begin{enumerate}
 \item La coincidencia entre las instrucciones de la clase compartida
 y la clase del notebook mediante sus \'arboles de sintaxis.
 \item La equivalencia de las contracciones \code{einsum} con sumas
 expl\'icitas sobre puntos de cuadratura.
 \item La suma de los dos bloques elementales que contribuyen
 a \code{G[0:2,4:6]}.
 \item La reconstrucci\'on del bloque diagonal $H^s_{00}$.
 \item Que \code{matrices} no modifica \code{self.Hs}.
 \item La coincidencia entre resolver manualmente con $H,G$ y
 utilizar \code{bem.resolver(40.0)}.
 \item Que duplicar $p$ duplica el desplazamiento complejo, como exige
 la linealidad.
\end{enumerate}

La suma de bloques de cada fila de $H^s$ es cero hasta errores
aproximadamente de $1.18\times10^{-16}$ en magnitud m\'axima.
El residuo hacia atr\'as de la soluci\'on es
\begin{equation}
 \eta=
 \frac{\norm{H\vect u-G\vect t}_2}
 {\norm H_F\norm{\vect u}_2+\norm{G\vect t}_2}
 \simeq3.30\times10^{-17}.
\end{equation}
Un residuo peque\~no verifica la soluci\'on del sistema ensamblado;
no elimina el error de malla o cuadratura.

\subsection{Lo que la clase presupone}
No es un programa gen\'erico que valide cualquier entrada.
Supone radio positivo, malla v\'alida, material compatible y funciones
auxiliares disponibles. El modelo es un c\'irculo centrado en el origen
con presi\'on conocida en toda su frontera.

A frecuencia cero, los n\'ucleos din\'amicos no se eval\'uan directamente
de esta manera; el problema est\'atico libre tiene movimientos r\'igidos.
Adem\'as, la resta de n\'ucleos cerca de la fuente puede perder precisi\'on.
No se debe extrapolar el ejemplo comprobado a cualquier frecuencia,
geometr\'ia o refinamiento.

\subsection{Qu\'e cambios requerir\'ian reconsiderar el c\'odigo}
\begin{itemize}
 \item \textbf{Otra geometr\'ia:} cambiar malla y c\'alculo de normales
 nodales; \code{nodes/R} no es general.
 \item \textbf{Otro sentido de recorrido:} revisar el signo de las normales.
 \item \textbf{Condiciones mixtas:} reorganizar las columnas conocidas y
 desconocidas; no basta con usar el mismo vector de presi\'on.
 \item \textbf{Materiales por zonas:} formular dominios e interfaces.
 \item \textbf{Campo interior:} agregar una evaluaci\'on de la representaci\'on
 integral despu\'es de resolver la frontera.
\end{itemize}
Son posibles ampliaciones, no funcionalidades implementadas por
estas 49 l\'ineas.

\section{Preguntas para comprobar que se entendi\'o}
\subsection*{1. Por qu\'e \texttt{self.N} vale 16 y $H$ mide $32\times32$?}
Hay 16 nodos y dos componentes por nodo. Cada nodo de colocaci\'on
aporta dos ecuaciones y cada nodo de campo dos componentes desconocidas.

\subsection*{2. \texttt{loc[0]} vale lo mismo que \texttt{ic}?}
No. \code{ic} es un n\'umero de nodo global.
\code{loc[0]} es la posici\'on local 0, 1 o 2 dentro de la conectividad,
si existe una coincidencia.

\subsection*{3. Por qu\'e se usa \texttt{J[:,None]}?}
Para convertir el vector de longitudes en una columna que divide
o multiplica cada fila del arreglo correspondiente.

\subsection*{4. Qu\'e significa la letra \texttt{q} en \texttt{einsum}?}
Es el punto de cuadratura. Si desaparece despu\'es de la flecha,
se suma sobre todos esos puntos.

\subsection*{5. Por qu\'e se usa \texttt{+=} en el ensamblaje?}
Porque un bloque global puede recibir contribuciones de distintos
elementos. Una asignaci\'on simple borrar\'ia contribuciones anteriores.

\subsection*{6. La diagonal de $H^s$ se pone igual a $I/2$?}
No. Se obtiene como menos la suma de los dem\'as bloques de su fila.
Incluye t\'ermino libre y contribuci\'on integral est\'atica.

\subsection*{7. Por qu\'e en la parte din\'amica no se omite la diagonal?}
Porque all\'i se integra $T-T^s$, cuya singularidad principal ya fue
retirada. La parte est\'atica diagonal se incorpor\'o por separado.

\subsection*{8. \texttt{reshape(-1,2)} calcula desplazamientos nuevos?}
No. Reorganiza el vector ya resuelto en filas de dos componentes.

\subsection*{9. \texttt{resolver} devuelve una amplitud \'unica?}
No. Devuelve un valor radial complejo por nodo.
Tomar magnitudes o promediarlas son operaciones posteriores.

\subsection*{10. Qu\'e se repite al pedir otra frecuencia?}
Se eval\'uan de nuevo los n\'ucleos din\'amicos, se ensamblan las
matrices y se resuelve. La geometr\'ia de cuadratura y $H^s$ se reutilizan.

\section{C\'omo utilizar y compilar este documento}
Este archivo \LaTeX{} es aut\'onomo: contiene el listado, los esquemas
TikZ y los n\'umeros del ejemplo. No necesita cargar figuras de la gu\'ia
anterior ni acceder al notebook para compilarse.

Desde su carpeta:
\begin{lstlisting}[language=bash]
latexmk -pdf -interaction=nonstopmode -halt-on-error \
  -outdir=.build explicacion_bem_cilindro.tex
cp .build/explicacion_bem_cilindro.pdf explicacion_bem_cilindro.pdf
\end{lstlisting}
Se requiere una distribuci\'on \LaTeX{} con espa\~nol, TikZ,
\code{listings}, \code{tcolorbox} y \code{latexmk}.
La numeraci\'on del listado y las referencias se actualizan
autom\'aticamente al compilar.

\begin{idea}[La clase en una sola frase]
El constructor prepara \emph{d\'onde y c\'omo integrar};
\code{matrices(w)} integra \emph{las influencias de esa frecuencia};
\code{resolver(w,p)} aplica \emph{la carga conocida},
resuelve \emph{los desplazamientos desconocidos}
y devuelve \emph{su componente radial}.
\end{idea}

\end{document}
